% book/chapters/ch04-agents.tex
\chapter{Agents and Subagents}
\label{ch:agents}

\index{agents}\index{subagents}

Claude Code can delegate subtasks to \emph{subagents} — specialized assistants
that run in their own isolated context windows. A subagent explores files,
reviews code, or executes a multi-step procedure without consuming tokens in
your main conversation. When it finishes, it returns a summary and the work is
done.

Several built-in subagents ship with Claude Code; you can also define custom
ones for your own domain.

\section{Why Use Subagents?}
\label{sec:agent-why}

\begin{itemize}
  \item \textbf{Context isolation.} A codebase search can read hundreds of
        files. Running it in a subagent keeps those token-heavy reads out of
        your main session.
  \item \textbf{Tool restriction.} A read-only subagent cannot accidentally
        modify files.
  \item \textbf{Cost optimization.} Lightweight tasks can be routed to a
        faster, cheaper model (e.g.\ Haiku).
  \item \textbf{Reuse.} A custom subagent defined once is available across
        all projects.
\end{itemize}

\section{Built-In Subagents}
\label{sec:built-in-agents}

\index{Explore agent}\index{Plan agent}\index{General-purpose agent}

\begin{table}[ht]
\centering
\caption{Built-in subagents in Claude Code}
\label{tab:built-in-agents}
\begin{tabularx}{\textwidth}{llXl}
\toprule
\textbf{Subagent} & \textbf{Model} & \textbf{Tools} & \textbf{Purpose} \\
\midrule
\textbf{Explore}          & Haiku            & Read-only
  & Fast file search and codebase understanding \\
\textbf{Plan}             & Inherits parent  & Read-only
  & Context gathering during plan mode \\
\textbf{General-purpose}  & Inherits parent  & All tools
  & Complex multi-step tasks requiring read and write \\
statusline-setup          & Sonnet           & Read, Edit
  & Configuring the status line (\texttt{/statusline}) \\
claude-code-guide         & Default          & Read, Bash, web search
  & Answering questions about Claude Code features \\
\bottomrule
\end{tabularx}
\end{table}

\begin{notebox}
A subagent can delegate to other subagents only if its \texttt{tools}
field explicitly includes the \texttt{Agent} tool (with or without an
allowlist of permitted children, e.g.\
\texttt{tools: Agent(researcher, worker), Read}). When \texttt{Agent} is
omitted from the \texttt{tools} field, the subagent cannot spawn other
subagents.
\end{notebox}

Claude automatically routes tasks to the appropriate subagent. You can also
delegate explicitly:

\begin{lstlisting}[caption={Explicit subagent delegation}]
"Use a subagent to review the auth module for edge cases"
"Use a subagent to investigate how token refresh is handled"
\end{lstlisting}

\section{Creating Custom Subagents}
\label{sec:custom-agents}

\index{agents!custom}

Custom subagents are defined as Markdown files in \texttt{.claude/agents/}
(project scope) or \texttt{\textasciitilde/.claude/agents/} (personal scope).

\subsection{File Format}

\begin{lstlisting}[caption={Custom subagent definition: \texttt{.claude/agents/security-reviewer.md}}]
---
name: security-reviewer
description: Reviews code for security vulnerabilities. Use when asked to
             audit code, check for injection flaws, or review auth logic.
tools: Read, Grep, Glob, Bash
model: opus
---
You are a senior security engineer. Review code for:
- Injection vulnerabilities (SQL, XSS, command injection)
- Authentication and authorization flaws
- Hardcoded secrets or credentials
- Insecure data handling and storage

Provide specific line references and concrete suggested fixes.
Do not suggest changes unrelated to security.
\end{lstlisting}

\subsection{Frontmatter Fields}

\begin{table}[ht]
\centering
\caption{Custom subagent frontmatter fields}
\label{tab:agent-frontmatter}
\begin{tabularx}{\textwidth}{llX}
\toprule
\textbf{Field} & \textbf{Required} & \textbf{Description} \\
\midrule
\texttt{name}             & Yes & Identifier used for delegation \\
\texttt{description}      & Yes & Natural-language; Claude uses this to decide when to delegate \\
\texttt{tools}            & No  & Allowlist; omit to inherit all parent tools.
                                  Use \texttt{Agent(child1, child2)} to permit a specific set of sub-subagents \\
\texttt{disallowedTools}  & No  & Denylist; opposite of \texttt{tools} \\
\texttt{model}            & No  & \texttt{haiku}, \texttt{sonnet}, \texttt{opus} \\
\texttt{memory}           & No  & \texttt{user}, \texttt{project}, or \texttt{local} — enables persistent memory (see below) \\
\texttt{permissionMode}   & No  & Restrict the subagent's permission mode \\
\texttt{mcpServers}       & No  & Scope a server (e.g.\ Playwright) to a single subagent \\
\texttt{skills}           & No  & Skills to preload into the subagent \\
\texttt{maxTurns}         & No  & Hard cap on tool-call turns \\
\texttt{isolation}        & No  & \texttt{worktree} gives the subagent its own git worktree \\
\texttt{color}            & No  & Display color in agent listings \\
\bottomrule
\end{tabularx}
\end{table}

\subsection{Scope}

\begin{itemize}
  \item \texttt{.claude/agents/} — project scope; commit to share with the team
  \item \texttt{\textasciitilde/.claude/agents/} — personal scope; available
        across all your projects
\end{itemize}

\subsection{Persistent Subagent Memory}
\label{sec:agent-memory}

\index{subagents!persistent memory}

Set \texttt{memory: user}, \texttt{memory: project}, or
\texttt{memory: local} in a subagent's frontmatter and Claude Code grants
it a dedicated directory that survives across conversations:

\begin{table}[ht]
\centering
\caption{Subagent memory scopes}
\label{tab:agent-memory}
\begin{tabularx}{\textwidth}{llX}
\toprule
\textbf{Scope} & \textbf{Directory} & \textbf{Use} \\
\midrule
\texttt{user}    & \texttt{\textasciitilde/.claude/agent-memory/<name>/}
                 & Learnings that apply across all projects \\
\texttt{project} & \texttt{.claude/agent-memory/<name>/}
                 & Project-specific; commit to share with the team \\
\texttt{local}   & \texttt{.claude/agent-memory-local/<name>/}
                 & Project-specific but \texttt{.gitignore}'d \\
\bottomrule
\end{tabularx}
\end{table}

When memory is enabled, the subagent's system prompt is extended with
read/write instructions for its memory directory, and the first 200 lines
(or 25\,KB) of \texttt{MEMORY.md} are loaded on every invocation. The
subagent's Read, Write, and Edit tools are enabled automatically so it can
manage its own files. Memory is the subagent equivalent of the
\emph{auto memory} feature from \cref{ch:context} — the same mechanism,
scoped to one specialist.

\section{Forking the Current Conversation}
\label{sec:agent-fork}

\index{agents!fork}

You can fork the current session into a subagent that inherits full context so
far, then let it explore a risky or exploratory path without affecting your
main conversation:

\begin{lstlisting}[caption={Forking a session for an exploratory investigation}]
"Fork the conversation and use a subagent to experiment with
 the cache invalidation approach; report back what works"
\end{lstlisting}

\section{Background Agents and Agent Teams}
\label{sec:background-agents}

\index{background agents}\index{agent teams}

For workloads larger than a single session:

\begin{itemize}
  \item \textbf{Background agents} (\texttt{/agent-view}) run multiple
        Claude Code sessions in parallel, each in its own git worktree, and
        let you monitor them from one screen.
  \item \textbf{Agent teams} coordinate multiple sessions with shared tasks,
        messaging, and a designated lead agent. They can loop autonomously
        across many tasks while you spot-check recorded findings.
\end{itemize}

Both are covered in the official documentation at
\url{https://code.claude.com/docs/en/agent-teams}.

\section{Exercises}
\label{sec:agent-exercises}

\begin{exercisebox}[Exercises — define and use a custom subagent]
These exercises continue the running project.

\begin{enumerate}
  \item Create \texttt{.claude/agents/security-reviewer.md} using the
        frontmatter template from this chapter. Set
        \texttt{tools: Read, Grep, Glob, Bash} so the subagent cannot
        modify files. Commit it.
  \item In a session, delegate explicitly:
        \emph{``Use the security-reviewer subagent to scan
        \texttt{src/} for hardcoded secrets and unsafe shell calls.''}
        Confirm the response is a list of findings, and check that the
        subagent did \emph{not} write anything to disk (the tool restriction
        worked).
  \item Now \emph{remove} the explicit ``Use the security-reviewer
        subagent'' phrasing and just ask: \emph{``Audit \texttt{src/}
        for security issues.''} The \texttt{description} field of your
        agent should match, so Claude routes to it automatically. If it
        doesn't, refine the \texttt{description} until it does.
  \item Add a second custom subagent: \texttt{test-runner.md} with
        \texttt{model: haiku} and tools \texttt{Bash, Read}. Its job is
        to run the project's test suite and report failures verbatim.
        Verify that a question like ``Run the tests and show me what
        failed'' triggers the haiku-backed agent rather than the main
        Sonnet session.
\end{enumerate}
\end{exercisebox}

%% ---------------------------------------------------------------
\begin{summarybox}
\begin{itemize}
  \item Subagents run in their own context windows, preserving tokens in your
        main conversation and enabling tool restriction and model routing.
  \item Five built-in subagents: \textbf{Explore} (Haiku, read-only),
        \textbf{Plan} (read-only), \textbf{General-purpose} (all tools),
        plus the more specialized \texttt{statusline-setup} and
        \texttt{claude-code-guide}.
  \item Custom subagents are Markdown files in \texttt{.claude/agents/} with
        YAML frontmatter. \texttt{name} and \texttt{description} are
        required; \texttt{tools}, \texttt{model}, \texttt{memory},
        \texttt{permissionMode}, \texttt{mcpServers}, \texttt{skills},
        \texttt{maxTurns}, and \texttt{isolation} are optional.
  \item The \texttt{description} field is how Claude decides when to delegate
        automatically — write it clearly.
  \item Set \texttt{memory: user|project|local} in the frontmatter to give
        a subagent a persistent notes directory that survives across
        conversations.
  \item A subagent can spawn other subagents only when its \texttt{tools}
        allowlist includes \texttt{Agent} (optionally with a child allowlist).
        For parallel multi-session workflows, use background agents or
        agent teams.
\end{itemize}
\end{summarybox}

\nocite{*}
\printbibliography[heading=subbibliography,title={References}]
